%*******************************************************
% Abstract
%*******************************************************
\pdfbookmark[0]{Abstract}{Abstract}
\chapter*{Abstract}

Black-Box-Modelle bieten zwar oft eine hohe Vorhersagegüte, sind aber schwer zu interpretieren.
In diesem Zusammenhang stellt diese Projektarbeit das interaktive Tool \textit{Visual Explainable AI} vor, das Studierenden Verfahren des Fachgebiets Explainable Artificial Intelligence (XAI) auf einer visuellen Nutzeroberfläche bereitstellt, um die Ergebnisse von komplexen neuronalen Netzwerken verständlicher zu machen.
\textit{Visual Explainable AI} wurde mit den Best-Practices in der Softwaretechnik entwickelt und bietet in der vorliegenden Version Visualisierungen für die XAI-Methoden GRAD-CAM und LIME auf die vortrainierten neuronale Netzwerke ResNet und EfficientNetV2 an, lässt sich aber aufgrund der implementierten Software-Architektur Model-View-Presenter (MVP) mit Supervising Presenter um weitere neuronale Netzwerke und XAI-Verfahren leicht erweitern.
Zur Evaluierung der Wirksamkeit des Tools auf Studierende wurde eine Umfrage durchgeführt.
Die Ergebnisse der Umfrage zeigten u. A., dass in der Wahrnehmung der Befragten das Erkennen relevanter Eingabemerkmale sowie potenzieller Fehlvorhersagen und Grenzen von KI-Modellen durch das Tool unterstützt wurde, wobei das Vertrauen in die Modellentscheidungen trotz Nutzung des Tools überwiegend neutral blieb.
Weitere Untersuchungen mit einer größeren Stichprobe sind notwendig, um die Effektivität des Tools umfassender zu bewerten und weitere Möglichkeiten zur Verbesserung zu identifizieren.
